% -----------------------------
% УДИРТГАЛ
% -----------------------------

\addchaptertocentry{Удиртгал}
\chapter*{Удиртгал}


Орчин үеийн нийгэмд хүйсийн болон нийгмийн тэгш байдал, хүртээмжтэй орчин нь хүний эрх, боловсрол, хөдөлмөр эрхлэлт, байгууллагын соёл болон нийгмийн хөгжилтэй шууд холбоотой чухал асуудлын нэг болж байна. Хүн бүр ялгаатай байдлаасаа үл хамааран мэдээлэл авах, суралцах, ажиллах, нийгмийн амьдралд оролцох тэгш боломжтой байх нь хүртээмжтэй нийгмийг бүрдүүлэх үндэс юм.

Гэвч хүйсийн хэвшмэл ойлголт, ялгаварлан гадуурхалт, нийгмийн тэгш бус байдал, хүртээмжгүй орчин зэрэг асуудлууд өнөөг хүртэл тодорхой хэмжээнд оршсоор байна. Эдгээр асуудал нь хувь хүний эрх, боломжид нөлөөлөхөөс гадна байгууллага, боловсролын орчин болон нийгмийн хөгжилд сөрөг нөлөө үзүүлдэг. Иймээс иргэд хүйсийн болон нийгмийн тэгш байдал, харилцан хүндлэл, хүртээмжтэй орчны талаарх зөв ойлголттой байх шаардлагатай.

Энэ төрлийн мэдээлэл нь ихэвчлэн хууль эрх зүйн баримт бичиг, бодлогын баримт бичиг, гарын авлага, сургалтын материал зэрэг олон эх сурвалжид тархсан байдаг. Мөн зарим мэдээлэл нь урт, мэргэжлийн нэр томьёо ихтэй байдаг тул хэрэглэгч шаардлагатай мэдээллээ богино хугацаанд олж авах, ойлгох, хэрэглэхэд хүндрэлтэй байдаг.

Сүүлийн жилүүдэд байгалийн хэл боловсруулалт (NLP) болон том хэлний загвар (LLM)-т суурилсан хиймэл оюун ухааны технологи эрчимтэй хөгжиж, асуултад хариулах, мэдээлэл хүргэх, хэрэглэгчтэй харилцах чиглэлээр өргөн ашиглагдаж байна. Гэсэн хэдий ч зөвхөн том хэлний загварт тулгуурласан систем нь эх сурвалжгүй эсвэл буруу хариулт үүсгэх эрсдэлтэй байдаг. Иймээс мэдлэгийн сангаас холбогдох мэдээллийг хайж, баримтад тулгуурлан хариулт үүсгэдэг RAG (Retrieval-Augmented Generation) архитектур нь илүү найдвартай шийдэл юм.

Иймд энэхүү дипломын ажил нь хүйсийн болон нийгмийн тэгш, хүртээмжтэй байдлын талаарх мэдээллийг хэрэглэгчдэд ойлгомжтой, хүртээмжтэй байдлаар хүргэх RAG суурьтай хиймэл оюун ухаант чатбот системийг судалж, боловсруулахад чиглэж байна.

\section*{Зорилго}

Энэхүү дипломын ажлын зорилго нь хэрэглэгчийн асуултад хүйсийн болон нийгмийн тэгш байдал, хүртээмжтэй орчны талаарх мэдээллийг баримтад тулгуурлан тайлбарлах, мэдлэгийн сангаас холбогдох мэдээллийг хайж татах RAG (Retrieval-Augmented Generation) суурьтай хиймэл оюун ухаанд суурилсан чатбот системийг боловсруулахад оршино.

\section*{Зорилтууд}

Дипломын ажлын зорилгыг хэрэгжүүлэхийн тулд дараах зорилтуудыг дэвшүүлэв.

\begin{enumerate}

\item Хүйсийн болон нийгмийн тэгш байдал, хүртээмжтэй байдлын онолын үндсийг судлах

\item AI чатбот, байгалийн хэл боловсруулалт (NLP) болон том хэлний загвар (LLM) технологийн талаар судалгаа хийх

\item Баримтад тулгуурласан хариулт үүсгэх RAG архитектурын зарчмыг судлах

\item Чатбот системийн архитектур болон мэдлэгийн сангийн бүтцийг боловсруулах

\item Хэрэглэгчийн асуултад хариулах чатбот системийн прототипийг хэрэгжүүлэх

\item Системийн ажиллагааг туршилтын асуултын багц ашиглан үнэлэх

\end{enumerate}